


\begin{figure}[h]
\centering
\vspace{-15pt}
\includegraphics[width=0.9\linewidth]{figures/sotacomp_plot2.pdf}
\includegraphics[width=\linewidth]{figures/teaser_v3.png}
\vspace{-20pt}
\caption{
    Top: Comparison of \modelname~to state-of-the-art text-to-motion models (DART \cite{DartControl}, LoM \cite{chen2024language}, GoToZero \cite{Fan_2025_ICCV}, and MoMask \cite{momask}). Bottom: Example results generated by \modelname~(retargeted to different characters). 
}
\label{fig:teasercomp}
\end{figure}


\section{Introduction}
\label{sec:intro}





The creation of high-quality 3D content has become a cornerstone of modern digital experiences, driving innovation in fields ranging from virtual and augmented reality (VR/AR) and interactive entertainment to robotics and digital human applications. 
Traditionally, generating realistic 3D human animations has required complex technical skills and a significant amount of time for well-trained animators, even with the assistance of expensive motion capture (MoCap) systems.
In recent years, the advance of generative AI has opened a new frontier: text-to-motion generation, a paradigm that aims to democratize 3D human animation by allowing creators to generate complex motions simply by describing them in natural language.

Significant progress in this domain has been fueled by various motion generation models \cite{mdm,momask,DartControl,motionlcm,remodiffuse,motiondiffuse,mld,Kim2025,Guo_2025_ICCV,Xiao_2025_ICCV,Pinyoanuntapong2025MaskControl} and motion datasets \cite{Guo2022,mahmood2019amass,Mason2022,Li2023,Kim2025,Harvey2020,Kobayashi2023,Guo2025}, with recent efforts exploring scaling motion generation with Large Language Models (LLMs) \cite{t2mgpt,motiongpt_nips,chen2024language,Lu_2025_CVPR,cao2025real,Fan_2025_ICCV}. 
On one hand, models such as MoMask \cite{momask} and DART \cite{DartControl} have shown promising results. However, being relatively small in model size, their generative capabilities are often limited, struggling to understand user instructions and produce motions of high complexity. 
On the other hand, models like LoM \cite{chen2024language} and GoToZero \cite{Fan_2025_ICCV}, have leveraged Large Language Models to generate motions by expanding vocabulary with discrete motion tokenizers. 
While this approach enhances semantic diversity, the quantization process inherent in discrete representations often leads to a degradation in motion quality, resulting in outputs that lack smoothness and naturalness. 
These divergent approaches highlight several persistent challenges that hinder the widespread adoption and commercial viability of text-to-motion models. First, motion fidelity and realism remain a primary concern; models frequently produce results with noticeable artifacts such as foot sliding, unnatural physics, or the aforementioned lack of smoothness. Second, the semantic alignment between the input text and the output motion remains a significant hurdle, as models often fail to accurately interpret complex or nuanced instructions. Third, unlike in image and video generation, the scaling of motion generation diffusion models has been relatively underexplored, leaving their full potential for enhanced quality and instruction-following largely untapped. Finally, all these challenges are compounded by a data bottleneck: the lack of large-scale, diverse, and meticulously cleaned and annotated motion datasets.


To address these critical gaps, we present HY-Motion 1.0, a new series of state-of-the-art, large-scale models for text-to-motion generation. 
Our work marks a significant leap forward by being the first to successfully scale a Diffusion Transformer (DiT)-based architecture \cite{DIT}, which leverages a flow matching objective \cite{Lipman2022,Esser2024,labs2025flux1kontextflowmatching}, to the billion-parameter scale in the motion generation domain. 
This scaling, combined with a novel and comprehensive training methodology, enables HY-Motion 1.0 to achieve unprecedented levels of motion quality and instruction-following precision, significantly outperforming existing open-source solutions.
To the best of our knowledge, we are the first to successfully scale a Diffusion Transformer (DiT)-based architecture \cite{DIT} to a billion-parameter scale, outperforming previous open-source state-of-the-art methods by a substantial margin, in terms of motion quality and the instruction-following capability.

Specifically, the foundation of our success lies in a full-stage training paradigm. 
Firstly, a large-scale pretraining is conducted on a motion dataset over 3,000 hours, including various motion categories and variations featuring maximal diversity and coverage, to grasp the motion structure and the semantic correlation between text and motion.
Secondly, we conduct fine-tuning on 400-hour high-quality text-motion pairs curated through meticulous motion filtering and caption filtering, resulting in better motion quality and more accurate instruction-following.
Crucially, we introduce a final alignment stage using reinforcement learning (RL) from both human feedback and reward models to ensure the generated motions are not only physically plausible but also precisely aligned with human preferences and textual commands. 
This entire framework is supported by our newly developed meticulous data processing pipeline, which integrates automatic motion data cleaning, processing, and captioning with extensive manual refinement, resulting in a dataset with the industry's widest coverage of over 200 distinct motion categories.

In this paper, we detail the data pipeline, model architecture, and training process of HY-Motion 1.0. Our primary contributions can be summarized as follows:

\begin{itemize}
    \item \textbf{Scaling Law for Text-to-Motion}: We are the first to successfully scale a DiT-based flow matching model to over one billion parameters for text-to-motion generation, demonstrating significant improvements in instruction following capabilities of text-to-motion models.
    \item \textbf{Comprehensive Full-Stage Training Paradigm}: We implement a three-stage training framework for the text-to-motion model training — large-scale pretraining, high-quality fine-tuning, and reinforcement learning — that holistically enhances model performance.
    \item \textbf{Meticulous Data Curation Pipeline}: We designed and implemented a meticulous data curation pipeline, including both automated processing and extensive manual refinement for data cleaning and captioning. This rigorous process has yielded a large-scale, high-quality motion dataset, 
    enabling the training of a large-scale network capable of generating high-quality motions accurately following the text instruction.
    \item \textbf{Open-Source Release}: We are releasing the HY-Motion 1.0 models to the open-source community to foster further research and accelerate the development of commercially viable 3D motion generation technologies.
\end{itemize}







